\begin{center}
    {\LARGE \textbf{2011无机化学}} \\[2ex]
    {\large 时量180分钟 \quad 满分100分}
\end{center}

\vspace{0.5cm}
\noindent\textbf{注意事项:}
\begin{enumerate}[label=\arabic*., parsep=0pt, itemsep=0pt]
    \item 答题前,考生先将自己的姓名、学号、考场号、座位号填写在答题卡上,将条形码准确粘贴在考生信息条形码粘贴区。
    \item 考试结束后,将本试卷与答题纸一并收回。
    \item 回答各个问题时,将答案写在答题纸上,写在本试卷上无效。
\end{enumerate}

\vspace{0.5cm}
\noindent\textbf{一、完成并配平下列反应的方程式（24 pts.）}

\begin{enumerate}[label=\arabic*., start=1, itemsep=1.5ex]
    \item \ce{(NH4)2Cr2O7}受热分解。
    \item 将\ce{KMnO4}和浓\ce{HCl}混合。
    \item 用\ce{NH2OH}还原\ce{AgBr}。
    \item 向\ce{SnCl2}溶液中加入\ce{FeCl3}溶液。
    \item 金属铜溶于氰化钠的水溶液中。
    \item 硝酸银溶液中滴加少量的硫代硫酸钠溶液。
\end{enumerate}

\vspace{0.5cm}
\noindent\textbf{二、解释下列实验现象（20 pts.）}

\begin{enumerate}[label=\arabic*., start=1, itemsep=1.5ex]
    \item 在煤气灯上加热\ce{KNO3}晶体时没有棕色气体生成，若\ce{KNO3}晶体中混有\ce{MgSO4}时则有棕色气体生成。
    \item 向\ce{K2Cr2O7}溶液中加入\ce{AgNO3}溶液，有砖红色沉淀生成；若再加入\ce{NaCl}溶液并煮沸，沉淀变成白色。
    \item 将\ce{MnSO4}、\ce{(NH4)2S2O8}和\ce{H2SO4}溶液混合后水浴加热，有棕褐色沉淀生成；若加入少量\ce{AgNO3}溶液，棕褐色沉淀很快消失，溶液变红。
    \item 有人做实验时，发现\ce{SnS}能被\ce{Na2S}溶解，分析可能原因并设计实验加以证明。
\end{enumerate}

\vspace{0.5cm}
\noindent\textbf{三、回答下列各题（56 pts.）}

\begin{enumerate}[label=\arabic*., start=1, itemsep=1.5ex]
    \item 写出具有下列电子构型的元素的名称、符号在周期表中的位置:
    \begin{enumerate}[label=(\arabic*), itemsep=1ex]
        \item $3d^74s^2$
        \item $4d^45s^1$
        \item $5s^25p^1$
        \item $5d^96s^1$
    \end{enumerate}
    \item 给出下列分子或离子的几何构型并指出中心原子的杂化方式。
    \begin{enumerate}[label=(\arabic*), itemsep=1ex]
        \item \ce{COCl2}
        \item \ce{[Ni(CN)4]^{2-}}
        \item \ce{ICl3}
        \item \ce{XeF2}
        \item \ce{[Zn(NH3)4]^{2+}}
    \end{enumerate}
    \item 写出\ce{N2}，\ce{N2^+}，\ce{O2}，\ce{O2^+}的分子轨道式，计算键级；并比较\ce{N2}和\ce{N2^+}；\ce{O2}和\ce{O2^+}的稳定性。
    \item 比较下列各组配位化合物的稳定性，并说明原因。
    \begin{enumerate}[label=(\arabic*), itemsep=1ex]
        \item \ce{[Cu(NH3)4]^{2+}}与\ce{[Cu(en)2]^{2+}}
        \item \ce{[Ag(S2O3)2]^{3-}}与\ce{[Ag(NH3)2]^+}
        \item \ce{[Fe(CN)6]^{4-}}与\ce{[Fe(CN)6]^{3-}}
        \item \ce{[Cu(NH3)2]^+}与\ce{[Ag(NH3)2]^+}
    \end{enumerate}
    \item 已知的热分解温度数据如下：
    
    \vspace{1ex}
    \begin{tabular}{l r @{\qquad} l r @{\qquad} l r}
        \ce{MgSO4} & 895$^\circ$C &
        \ce{CaSO4} & 1149$^\circ$C &
        \ce{SrSO4} & 1374$^\circ$C \\

                   &              &
        \ce{CdSO4} & 826$^\circ$C &
        \ce{MnSO4} & 755$^\circ$C \\

                   &              &
                   &              &
        \ce{Mn2(SO4)3} & 200$^\circ$C \\
    \end{tabular}
    \vspace{1ex}
    \begin{enumerate}[label=(\arabic*), itemsep=1ex]
        \item 试分析影响硫酸盐热分解温度的因素；
        \item 与硫酸盐相比，碳酸盐和硅酸盐的热稳定性如何，为什么？
    \end{enumerate}
    \item 比较\ce{HClO}、\ce{HClO3}、\ce{HClO4}酸性及氧化性强弱，说明原因。
    \item 已知\ce{Co^{3+}}的$P=17800\,\mathrm{cm}^{-1}$。\ce{Co^{3+}}与下列配体形成配离子的$\Delta_0$为
    \[
    \begin{array}{lccc}
     & \ce{F^-} & \ce{H2O} & \ce{NH3} \\
    \Delta_0/\mathrm{cm}^{-1} & 13000 & 18660 & 23000
    \end{array}
    \]
    \begin{enumerate}[label=(\arabic*), itemsep=1ex]
        \item \ce{Co^{3+}}与这些配体形成的配位化合物中$d$电子的排布情况如何，是高自旋还是低自旋？
        \item 计算这些配位化合物的CFSE。
    \end{enumerate}
    \item 已知：$E^\ominus(\ce{Fe^{3+}/Fe^{2+}})=0.77V$，$E^\ominus\left(\ce{[Fe(SCN)5]^{2-}/Fe^{2+}}\right)=0.39V$
    求反应\ce{Fe^{3+} + 5SCN^- = [Fe(SCN)5]^{2-}}的平衡常数$K^\ominus$
\end{enumerate}

\vspace{0.5cm}
\noindent\textbf{四、制备与分离（20 pts.）}

\begin{enumerate}[label=\arabic*., start=1, itemsep=1.5ex]
    \item 用反应方程式表示下列转化过程：\ce{Co -> CoSO4 * 7H2O -> CoCl * 6H2O -> [Co(NH3)6]Cl3}
    \item 以软锰矿（\ce{MnO2}）为原料制备高锰酸钾，写出相关的实验步骤和各步反应方程式。
    \item 设计方案分离含有\ce{Ag^+}、\ce{Al^{3+}}、\ce{Co^{2+}}、\ce{Cr^{3+}}、\ce{Fe^{3+}}的混合溶液。
\end{enumerate}

\vspace{0.5cm}
\noindent\textbf{五、鉴别与推断（30 pts.）}

\begin{enumerate}[label=\arabic*., start=1, itemsep=1.5ex]
    \item \ce{PbO4}中的\ce{Pb}有两种不同的价态，设计实验加以验证。
    \item 现有下列六种磷的含氧酸盐固体。
    \ce{Na4P2O7}、\ce{NaPO3}、\ce{Na2HPO4}、\ce{NaH2PO4}、\ce{NaH2PO2}、\ce{NaH2PO3}
    请加以鉴定，写出相关的反应式。
    \item 有一固体混合物可能含有\ce{FeCl3}、\ce{NaNO2}、\ce{Ca(OH)2}、\ce{AgNO3}、\ce{CuCl2}、\ce{NaF}、\ce{NH4Cl}七种物质中的若干种。若将此混合物加水后，可得白色沉淀和无色溶液，在此无色溶液中加入\ce{KSCN}没有变化；无色溶液可使酸化后\ce{KMnO4}溶液紫色褪去；将无色溶液加热有气体放出。另外，白色沉淀可溶于\ce{NH3\cdot H2O}中。
    试根据上述现象，判断在混合物中，哪些物质肯定存在？哪些肯定不存在？哪些可能存在？简述理由，并用方程式表示。
    \item 将紫色晶体\ce{A}和蓝绿色晶体\ce{B}混合后溶于稀\ce{H2SO4}中，向其中加\ce{KOH}溶液得深棕色混合沉淀\ce{C}；向\ce{C}中加稀\ce{H2SO4}，沉淀部分溶解得黄棕色溶液\ce{D}；向\ce{D}中加过量\ce{NH4F}溶液得无色溶液\ce{E}。将不于稀\ce{H2SO4}的沉淀与\ce{KOH}，\ce{KClO3}固体共熔得绿色物质\ce{F}；将\ce{F}溶于水通入\ce{Cl2}，蒸发结晶又得化合物\ce{A}
    试给出\ce{A}、\ce{B}、\ce{C}、\ce{D}、\ce{E}、\ce{F}的化学式，并写出生成\ce{E}及生成\ce{F}的化学反应方程式。
\end{enumerate}